\chapter{估计样本量：应用与实例}
\chapterauthors{Warren S. Browner、Thomas B. Newman、Mark J. Pletcher}
\sourcepage{126}
第 5 章介绍了样本量估计的基本原理，本章给出几种照步骤操作的方法。先讨论分析性研究或实验，包括多变量分析等特殊问题；再介绍主要为描述性的研究；随后讨论样本量固定时如何处理、怎样提高效能，以及资料似乎不足时如何估算。最后列出常见错误。附录提供几种基本方法的查阅表。

我们始终教初学者自己估算样本量，可以使用本章表格、统计软件，或为加利福尼亚大学旧金山分校（UCSF）临床研究培训项目开发的 \url{https://www.sample-size.net}。即使设计需要比本章更复杂的方法，或资助方要求生物统计支持，在请教统计学家前，先作粗略估算也值得。你通常会惊喜地发现，它与专业计算相当接近！

\section{分析性研究与实验的样本量方法}
方法虽有不同，但都有以下步骤：
\sourcepage{127}
\begin{enumerate}
\item 陈述零假设及单侧或双侧备择假设。
\item 根据假设中预测及结局变量的类型，从表 6.1 选择适当检验。
\item 选择合理效应量，必要时估计测量变异程度。
\item 设定 alpha 和 beta。除非备择假设明确为单侧，否则采用双侧 alpha。
\item 使用适当附表、在线计算器或统计软件估计人数。
\end{enumerate}
\begin{bookbox}{表 6.1\quad 估计样本量的简单统计检验\textsuperscript{a}}
\begin{tabularx}{\linewidth}{YYY}\toprule
 & \multicolumn{2}{c}{结局变量}\\
预测变量 & 二分类 & 连续\\\midrule
二分类 & 卡方检验\textsuperscript{b} & $t$ 检验\\
连续 & $t$ 检验 & 相关系数\\\bottomrule
\end{tabularx}
\footnotesize\textsuperscript{a}其他无序或有序分类变量，以及计划使用其他分析方法时，见“其他考虑与特殊问题”。\\
\textsuperscript{b}卡方检验始终为双侧；其单侧等价形式是 $Z$ 统计量。
\end{bookbox}
即使某些参数的准确取值不确定，也应及早估算。等得太久，可能突然发现必须从头换用参数，甚至重新设计整个研究，所以本书很早介绍这一主题。并非所有分析性研究都适合下述三大类；一些常见例外将在“其他考虑与特殊问题”介绍。

\sourcepage{128}
\subsection{\texorpdfstring{$t$}{t} 检验}
\term{$t$ 检验}（t test）也称 Student $t$ 检验，Student 是其创立者的笔名。它常用于判断两组连续变量均值是否显著不同。虽然假定两组变量分布近似正态钟形曲线，但只要人数不太少（不少于约 30～40 人），也无极端离群值，就可用来估计大多数连续变量的样本量。

它通常比较队列研究或实验中的连续结局，例如母亲接受两种不同抗抑郁药后新生儿的出生体重；也可比较病例对照研究中的连续预测变量，此时比较病例与对照的预测变量均值。

用 $t$ 检验比较均值的研究（例 6.1a、6.1b），需要：
\begin{enumerate}
\item 陈述零假设并明确备择假设侧数。
\item 用两组连续变量均值之差表示效应量 $E$。
\item 以预期标准差 $S$ 表示连续变量变异。
\item 计算\term{标准化效应量}（standardized effect size）$E/S$。
\item 设定 alpha 和 beta。
\end{enumerate}
\begin{bookbox}{例 6.1a\quad 队列研究中使用 $t$ 检验估计样本量}
\textbf{问题：}在单用沙丁胺醇基础上加用异丙托溴铵，能否改善哮喘控制？研究者计划随机试验，测量治疗两周后的第一秒用力呼气容积（forced expiratory volume in 1 second，FEV$_1$）。
\sourcepage{129}
既往研究报告，沙丁胺醇治疗者的平均 FEV$_1$ 为 2.0 L，标准差 0.5 L。她希望检出两组均值至少 10\% 的差异。在双侧 $\alpha=0.05$、效能 0.80 下，两组各需要多少人？

\textbf{解答：}计算样本量所需的各项参数如下。零假设为：哮喘患者治疗两周后，单用沙丁胺醇者与联合使用沙丁胺醇和异丙托溴铵者的平均 FEV$_1$ 相同。双侧备择假设为：两组治疗两周后的平均 FEV$_1$ 不同。
\begin{itemize}
\item 单药均值 2.0 L；联合治疗均值 2.2 L；效应量 $E=2.2-2.0=0.2$ L，即 2.0 L 的 10\%。
\item 标准差 $S=0.5$ L，标准化效应量 $E/S=0.2/0.5=0.40$。
\item 双侧 $\alpha=0.05$；$\beta=1-0.80=0.20$。
\end{itemize}
在表 6A 最左列找到 0.40，向右与双侧 alpha 0.05、beta 0.20 列相交，每组需 100 人完成研究。考虑退出，还应招募更多。如果不可行，可调整设计，或接受只能检出更大的效应量。例 6.8 的配对测量可能提供解决办法。
\end{bookbox}
\sourcepage{130}
\begin{bookbox}{例 6.1b\quad 病例对照研究中使用 $t$ 检验估计样本量}
\textbf{问题：}血浆双氢睾酮是否与精原细胞瘤（一种睾丸癌）风险有关？研究者计划比较新发病例及同一社区对照的浓度。男性平均浓度约 1.5 nmol/L，标准差 0.5 nmol/L；疾病罕见，可假定对照相同。若希望检出至少 20\% 均值差异，在双侧 $\alpha=0.05$、效能 0.80 下，每组需要多少人？

\textbf{解答：}计算样本量所需的各项参数如下。零假设为：精原细胞瘤病例与对照的平均血清双氢睾酮水平相同。双侧备择假设为：精原细胞瘤病例的平均血清双氢睾酮水平与对照不同。
\begin{itemize}
\item 对照均值 1.5 nmol/L；病例均值 1.8 nmol/L；$E=1.8-1.5=0.3$ nmol/L，即 1.5 的 20\%。
\item $S=0.5$ nmol/L；$E/S=0.3/0.5=0.6$。
\item 双侧 $\alpha=0.05$；$\beta=0.20$。
\end{itemize}
表 6A 的 0.60 行与相应 alpha、beta 列相交，需 45 例病例及 45 例对照。
\end{bookbox}
\sourcepage{131}
如第 5 章所述，效应量难以选择。对 $t$ 检验，除 alpha 和 beta，只需一个标准化效应量，但这个无单位数值与临床现实有些脱节。最好先指定两组预期均值，求差，再除以标准差。虽然计算只需均值差、不需两个均值本身，这样做可使选择更贴近实际。

标准差包含真实人群变异与测量误差，十分关键：在给定效应量下，它越大，所需样本越大。通常可从文献或专家获得估计；但使用新问卷或仪器时，也可能没有有意义资料，见后文“资料不足时如何估计”。

连续测量的变化值，例如研究期间体重变化，常可作为结局，因为其标准差通常小于变量本身，从而减少样本量，见“使用配对测量”及例 6.8。

确定 alpha、beta 及标准化效应量后，建议用网站或统计软件计算。组大小相等时，也可用表 6A：从最左列找到效应量，再向右寻找所选 alpha 和 beta 下的每组人数。浏览本章附表，有助于理解各种选择如何影响人数。
\sourcepage{132}
当人数超过约 30、效能 0.80、双侧 alpha 0.05 时，有一个便捷近似（1）：
\[
\text{每个等大组的样本量}=\frac{16}{(\text{标准化效应量})^2}.
\]
例 6.1a 得 $16/(0.4)^2=100$ 人／组，与查表及在线计算器一致。

\subsection{卡方检验}
\term{卡方检验}（chi-squared test，$\chi^2$）比较两组具有某个二分类结局或预测特征的比例。它始终双侧，单侧假设对应单侧 $Z$ 检验。

队列、横断面研究或实验中，效应量为两组预期结局比例 $P_0$ 和 $P_1$ 的差。例如，研究除草剂暴露与非霍奇金淋巴瘤，$P_0$ 是未暴露者发病比例，$P_1$ 是暴露者发病比例。变异是两比例的函数，无须另行规定。

病例对照研究略有不同：$P_1$ 为病例中预期暴露比例，$P_0$ 为对照中预期暴露比例。例如，研究纯素饮食能否预防结肠癌，分别是病例和对照中的纯素食者比例。变异同样不必另行指定。
\sourcepage{133}
对于将用卡方检验或 $Z$ 检验比较两个比例的研究（例 6.2a、6.2b），估计样本量需要完成三步：陈述零假设，决定备择假设采用单侧还是双侧；用两组待比较的比例 $P_0$、$P_1$ 估计效应量及其变异；设定 alpha 和 beta。
\begin{bookbox}{例 6.2a\quad 队列研究中使用卡方检验估计样本量}
\textbf{问题：}练太极者是否比慢跑者更少发生背痛？文献提示慢跑者两年风险约 30\%。研究者希望显示太极使绝对风险至少降低 10 个百分点。在双侧 $\alpha=0.05$、效能 0.80 下，需多少人才能判断太极组两年发病率为 20\% 或更低？

\textbf{解答：}零假设为：慢跑者与练太极者的背痛发病率相同。双侧备择假设为：两组的背痛发病率不同。慢跑组的预期发病比例 $P_1=0.30$，太极组的预期发病比例 $P_0=0.20$，所以效应量为 $|0.30-0.20|=0.10$。设定双侧 $\alpha=0.05$、$\beta=1-0.80=0.20$。在表 6B.1 中找到两比例中较小者 0.20 所在行及差值 0.10 所在列，取该单元格三个数中的中间数，得出每组需有 313 人完成研究。
\end{bookbox}
\sourcepage{134}
\begin{bookbox}{例 6.2b\quad 病例对照研究中使用卡方检验估计样本量}
\textbf{问题：}妊娠晚期吃寿司是否与罕见结局前置胎盘有关？希望在双侧 alpha 0.05、效能 0.80 下检出至少 2.5 的比值比，病例与对照等大。网上调查提示约 25\% 孕妇在妊娠晚期吃寿司。需要多少病例及无并发症妊娠对照？

\textbf{解答：}零假设为：前置胎盘病例与对照中，在妊娠晚期吃寿司者的比例相同。双侧备择假设为：这两个比例不同。因为结局罕见，对照中吃寿司的比例近似于所有孕妇中吃寿司的比例，故 $P_0=0.25$。按下文的比值比公式，病例中吃寿司的预期比例为：
\[
P_1=\frac{2.5\times0.25}{(1-0.25)+(2.5\times0.25)}\approx0.45.
\]
效应量为 $|0.45-0.25|=0.20$。设定双侧 $\alpha=0.05$、$\beta=1-0.80=0.20$。在表 6B.1 的 0.25 行、0.20 列读取中间数，得出需要 98 例病例及 98 例对照。
\end{bookbox}
附录 6B 给出等大组下各种 alpha、beta 及比例组合的人数。
\sourcepage{135}
先在表 6B.1 或 6B.2 最左列找到 $P_0$、$P_1$ 较小者；必要时分别四舍五入到最接近的 0.05 或 0.01，再向右按差值与 alpha、beta 查每组人数。

组大小不等时，附表不够用，需要计算。例如例 6.2b 中，对照比病例容易找，每例配多个对照可减少病例数。在网站“Proportions-Sample size”计算器输入 alpha 0.05、beta 0.20，将第一组病例占总样本比例 $q_1$ 设为 0.25，即病例:对照为 1:3。$P_1=0.45$、$P_0=0.25$ 时，仅需 63 例病例及 190 例对照。总人数虽较多，$63+190>98+98$，招募却可能更可行。

\subsubsection{由风险比或比值比估计 \texorpdfstring{$P_0$}{P0}、\texorpdfstring{$P_1$}{P1}}
用\term{风险比}或\term{相对危险度}（risk ratio or relative risk）考虑效应量，通常很有帮助。对队列、横断面研究或实验，$\mathrm{RR}=P_1/P_0$，很容易互相换算。例如，研究每天使用社交媒体超过五小时的青少年，抑郁风险是否为较少使用者的两倍。若后者风险 $P_0=0.06$，则前者 $P_1=0.12$。

病例对照研究不能直接计算风险比（见附录 9B），必须使用\term{比值比}（OR）。
\sourcepage{136}
其中 $P_0$ 是对照暴露比例，$P_1$ 是病例暴露比例：
\[
\mathrm{OR}=\frac{P_1(1-P_0)}{P_0(1-P_1)},\qquad
P_1=\frac{\mathrm{OR}\,P_0}{(1-P_0)+\mathrm{OR}\,P_0}.
\]
例如，对照中过度使用社交媒体者预计占 15\%，希望检出 OR 为 2，则
\[
P_1=\frac{2\times0.15}{(1-0.15)+(2\times0.15)}=\frac{0.3}{1.15}=0.26.
\]
\subsubsection{由总体风险估计两组风险}
研究者通常能指定想检出的风险比，却未必知道有预测因素者 $P_1$ 和无该因素者 $P_0$ 的风险。如果知道预定样本总体风险 $P$，就能按给定 RR 估计。两组等大时：
\[
P_0=\frac{2P}{\mathrm{RR}+1},\qquad P_1=\mathrm{RR}\,P_0.
\]
例如，前瞻性队列研究低于或等于中位数的日照暴露，是否使老年女性髋骨折风险增至 1.5 倍，故 RR 为 1.5；
\sourcepage{137}
预计随访中约 1\% 骨折，即 $P=0.01$。则 $P_0=2\times0.01/(1.5+1)=0.008$，$P_1=1.5\times0.008=0.012$。组等大时，可用 $P$ 应等于两风险均值来检查。

更常见的是组不等大，需加入未暴露组占总样本的比例 $q_0$：
\[
P_0=\frac{P}{\mathrm{RR}+q_0(1-\mathrm{RR})},\qquad P_1=\mathrm{RR}\,P_0.
\]
例如，处于日照最低十分位者占 10\%，则 $q_0=0.90$。$P_0=0.01/[1.5+0.90\times(-0.5)]=0.00952$，$P_1=1.5\times0.00952=0.0143$。此时 $P$ 应等于两风险的加权平均。

\subsection{相关系数}
\term{相关系数}（correlation coefficient，$r$）在临床研究样本量计算中不常用，但预测与结局均连续时可用。它衡量线性关联强度，范围 $-1$～$+1$；绝对值越接近 1 越强，越接近 0 越弱。某些人群成人身高与体重高度相关，$r\approx0.9$，但许多生物学关联小得多。负值表示一个增加、另一个减少，例如儿童血铅与智商。

某些领域，如行为医学，常用相关系数，但用它估算样本量的缺点是缺乏直观意义。其平方 $r^2$ 表示结局变异，即方差中可由预测变量线性关联解释的比例，反之亦然。
\sourcepage{138}
小的 $r$，例如 $\leq0.3$，在大样本可显著，却可能没有多少临床或科学意义，因为最多“解释” $0.3^2=9\%$ 方差。

另一个往往更好的方法，是把两个连续变量之一二分，例如以中位数为界，再用 $t$ 检验估计样本量；其优点是能把效应量表达为两组之间的差异。若仍选择相关系数法（例 6.3），应先陈述零假设，决定备择假设采用单侧还是双侧；再用希望检出的相关系数 $r$ 的绝对值估计效应量；最后设定 alpha 和 beta。变异是 $r$ 的函数，表中已包含。
\begin{bookbox}{例 6.3\quad 横断面研究中使用相关系数估计样本量}
\textbf{问题：}吸烟者尿可替宁水平（反映当前吸烟强度）是否与骨密度相关？既往研究发现，自报每日吸烟支数与骨密度（g/cm$^3$）有一定程度的相关，$r=-0.3$；研究者预计，尿可替宁水平与骨密度的相关程度至少与此相同。双侧 $\alpha=0.05$、$\beta=0.10$ 时，需要招募多少名吸烟者？

\textbf{解答：}零假设为：吸烟者的尿可替宁水平与骨密度没有相关性。备择假设为：吸烟者的尿可替宁水平与骨密度存在相关性。
\sourcepage{139}
效应量取相关系数的绝对值，即 $|-0.3|=0.3$。设定双侧 $\alpha=0.05$、$\beta=0.10$。在表 6C 中读取 $r=0.30$ 所在行与相应 alpha、beta 所在列，得出需要 113 名吸烟者。
\end{bookbox}
附录 6C 按 $r$ 行与 alpha、beta 列给出总样本量，适用于检验总体无关联，即 $r=0$。若想检验是否不同于非零值，例如 $r=0.4$，应使用网站计算器。

\section{其他考虑与特殊问题}
\subsection{退出研究者}
每个抽样单位必须有资料可分析；虽已入组，但结局无法确定者，例如退出研究者，不计入有效样本量。若预计有人无法随访，应估计比例并相应增加招募人数。例如预计失访 20\%，就需把样本量乘以 $1/(1-0.20)=1.25$。

\subsection{分类与计数变量}
有序分类有逻辑顺序，例如无、轻、中、重疼痛；无序分类如血型则没有。虽然数学上用 $t$ 检验处理有序变量未必适当，
\sourcepage{140}
实际中若类别较多（至少六个）、观察在各类分布良好，且取平均有意义，可以把它近似连续变量。

其他情况下，最好略改假设，把变量二分。例如，英语能力分为几乎不会、会一点、熟练、流利及母语，考察其与急诊等待时间的关系时，可按“会一点及以下”与“熟练及以上”估算。

计数变量也类似。虽然有正式方法（2），但若至少六个取值且分布良好，可用 $t$ 检验近似；否则可在中位数附近二分，用卡方检验估计。

\subsection{生存分析}
比较两组生存或不发生不良事件，如癌症复发，的时间，可采用\term{生存分析}（survival analysis）（3、4）。晚期乳腺癌至复发月数看似连续，却不宜用 $t$ 检验，因为研究结束仍存活且未复发者没有该结局时间。同样，也可能比较每 100 人年等单位的结局发生率。两者都可先估计某时间内两组发生结局的比例，用卡方检验近似估算。
\sourcepage{141}
若多数人最终都会发生结局，例如晚期胰腺癌死亡，更节省总人数的办法，是在约一半总体结局已发生的时点，估计两组比例。例如，比较标准与实验治疗后的无病生存，预计两年时标准组死亡 60\%、实验组 40\%，可用两年生存这一二分类结局计算。网站计算器可以给出更精细的生存分析估算。

\subsection{群组样本}
某些设计按群组抽取参与者（第 12 章）。例如，研究继续医学教育能否提高患者戒烟率，把 20 家诊所随机分到干预组，另 20 家作对照。一年后，每家随机检查 50 名研究开始时吸烟患者的病历。样本量是 40 家诊所，还是 2000 名患者？答案介于两者，取决于同诊所患者戒烟倾向的相似程度，相对全体患者的相似程度有多大。除非已有研究，常需小型研究估计。相关方法已有多种（5～7），但较难，通常需要统计学家帮助。

\subsection{匹配}
\sourcepage{142}
研究者可能因多种原因采用匹配设计（第 10 章）。本章忽略匹配的方法，通常仍可给出合理估计，除非病例对照中的暴露、或队列中的结局，与匹配变量高度相关。更准确估算需要指定匹配对暴露或结局的相关程度，可用标准方法（8）或软件。

\subsection{多变量调整与其他特殊分析}
观察性研究可能预期某些变量混杂预测与结局关联（第 10 章），计划分析时调整。如果主要假设检验包括调整，样本量也应考虑。

调整混杂通常会增加人数要求（9、10），增加程度取决于混杂因素的流行程度、预测与混杂因素关联强度，以及混杂因素与结局关联强度等。作用复杂，没有普遍规则。

线性及 logistic 回归等方法可调整混杂；Cox 比例风险分析还可调整随访长度差异，以风险率比衡量关联。使用这些方法时，有对应样本量方法（3、11～14）。潜在遗传危险因素或候选基因研究（15～17）、经济学研究（18～20）、剂量反应研究（21）及多于两组的研究（22），也有相应方法，
\sourcepage{143}
互联网可提供这些较复杂方法的资料。

至少对初学者，先假设使用卡方或 $t$ 等简单检验较容易，也可检查复杂方法结果。例如，病例对照研究出生体重这一连续变量与儿童脑瘤这一事件时间结局的关系，即使最终拟用 Cox 分析，也可先用 $t$ 检验粗估。简单方法通常与复杂方法相近。不过，申请高额资助时，应咨询有经验的生物统计学家。即使评审者知道结局风险及效应量靠猜测，仍会期待适当的复杂方法。统计学家参与估算，也表明有管理和分析数据所需的合作者。她还会在设计和实施中作出许多贡献，并会乐于与已思考问题、搜集参数资料并作过初估的临床研究者合作。

\subsection{等效性与非劣效性试验}
有时目标是排除预测与结局之间存在较大关联。\term{等效性试验}（equivalence trial）检验新药是否与已有药物疗效近似相同；希望效应量为零或很小，给样本量规划带来挑战。\term{非劣效性试验}（noninferiority trial）是其单侧形式，检验新药是否至少没有明显更差（第 12 章）。
\sourcepage{144}
计算较复杂（23～26），有经验统计学家的建议很有帮助。粗估可选择很小、无临床意义的效应量，如平均血糖相差 5 mg/dL，并要求较高效能，例如 0.90 或 0.95，来拒绝零假设。但小效应和高效能往往需要很大样本；两者相比，非劣效性单侧设计可使用较少人数或较低 alpha。

另一问题是，通常零假设范式所提供的保护不再适用。在比较有效药物与安慰剂的常规研究中，零假设范式有助于防止 I 类错误，即错误拒绝零假设：不精密测量或严重失访等设计、实施问题，会使零假设更难被拒绝。因此，想拒绝零假设的研究者有很强的动力把研究做好。但非劣效性研究的目标是找不到差别；设计和实施不良会模糊组间差异，使本来存在的差异更容易漏掉，这些保护也就不起作用。

\section{描述性研究的样本量方法}
包括诊断检查在内的描述性研究，原理有所不同。它不依赖显著性检验，因此不适用效能及零、备择假设概念，而是估计均值、比例等描述统计，样本规模决定估计精度。不过，描述性研究常兼问分析性问题，例如调查老年门诊患者抑郁患病率时，也想知道预测因素。
\sourcepage{145}
此时也应按分析性问题估算，避免后来发现更感兴趣的问题效能不足。

描述性估计常伴有\term{置信区间}（confidence interval），即样本均值或比例周围的取值范围，衡量样本估计精度。研究者指定 95\% 或 99\% 等置信水平。较高水平，例如 99\%，区间比 90\% 更宽，也更可能涵盖总体真实值。

区间宽度取决于样本量。例如，想估计接受网上替代课程的医学生美国执照考试平均分。50 人样本可能估计全体均值 215，95\% 区间 205～225；20 人的小样本可能均值相似，但通常较不精确，区间更宽。估算时，应指定置信水平及所需宽度，再查附录 6D、6E 或在线计算。

\subsection{连续变量}
连续变量常报告均值置信区间。需要估计变量标准差、指定精度即区间总宽，并选择置信水平，如 95\%、99\%。
\begin{bookbox}{例 6.4\quad 连续变量描述性研究的样本量}
\sourcepage{146}
\textbf{问题：}研究者希望估计某城市三年级学生的平均阅读年级水平，使 95\% 置信区间达到均值上下各 $0.25$ 个年级。类似城市的既往研究发现，阅读年级水平的标准差为 1.4 个年级。

\textbf{解答：}变量的标准差为 1.4 个年级；区间总宽为 0.5 个年级，即均值上下各 0.25 个年级。因此，标准化区间宽度为总宽除以标准差，即 $0.5/1.4=0.35$。置信水平设为 95\%。在表 6D 中读取标准化宽度 0.35 所在行与 95\% 置信水平所在列，得出需要 126 名三年级学生。
\end{bookbox}
使用附录 6D，应先以总宽除以标准差，再按标准化宽度和置信水平查总人数。

\subsection{二分类变量}
二分类描述性结果，可表示为某一取值比例周围的置信区间。诊断检查灵敏度和特异度乍看似连续，实际是比例（第 13 章），也属于此类。需要估计总体特征比例；若超过一半，则按无该特征比例估算。再指定总宽和置信水平。
\sourcepage{147}
附录 6E 按预期比例行及总宽、置信水平列给出人数。例 6.5 研究灵敏度，得到的是所需患病者人数；研究特异度则需计算未患病者人数。第 13 章其他参数，例如受试者工作特征（receiver operating characteristic，ROC）曲线、似然比及可靠性，也有特定方法（27～29）。
\begin{bookbox}{例 6.5\quad 二分类变量描述性研究的样本量}
\textbf{问题：}研究甲状腺癌新诊断检查的灵敏度。先导研究提示患病者 80\% 阳性，若希望 95\% 区间为 $0.80\pm0.05$，需多少患病者？

\textbf{解答：}因为预期灵敏度 0.80 大于 0.5，按预期假阴性结果的比例 0.20 估算样本量。区间总宽为 0.10，即预期比例上下各 0.05；置信水平设为 95\%。在表 6E 的 0.20 行、总宽 0.10 列读取中间数，它对应 95\% 置信水平，得出需要 246 名甲状腺癌患者。
\end{bookbox}

\sourcepage{148}
\section{样本量固定时怎么办}
研究规划常在初估人数无法获得、负担不起或两者兼有时面临艰难选择。下一步往往是：按现实人数估计可检出效应量，再决定值不值得研究。有经验者常直接从这里开始。

二次数据分析可能在设计前就已确定人数；即使从零设计，可用或可负担人数通常也有限。因此，研究者常从固定或现实人数“反推”在合理效能下可检出多大效应。这也是不应把样本量估计当成不可改变规定的原因。

反推时（例 6.6），通常估计给定效能，如 80\%，下可检出效应；较少见的是给定效应量求效能。可用附表，必要时插值，或用统计软件和网站计算。
\begin{bookbox}{例 6.6\quad 人数固定时估计可检出效应量}
\textbf{问题：}研究给双胞胎新生儿母亲六周冥想课程，与只发放放松宣传册相比，是否降低 0～30 分量表的应激。培训期间可接触 200 人，先导研究提示一半，即 100 人，可能愿意参加并完成。两组预期标准差均为 5 分，
\sourcepage{149}
双侧 alpha 0.05、beta 0.20 时，可检出多大组间差？

\textbf{解答：}在表 6A 中，沿双侧 $\alpha=0.05$、$\beta=0.20$ 所在列（中间一组三列的最右列）向下查找，可见每组 45 人能检出标准化效应量 0.6，对应应激评分相差 $0.6\times5=3$ 分。研究者实际约有每组 50 人，因此能检出略小于 3 分的组间差异。
\end{bookbox}
一般认为应有至少 80\% 效能检出合理效应，但 80\% 并无魔力：效能低也可能幸运地获得显著结果；50\% 效能仍有一半机会发现总体真实效应。因此，若成本很低，例如分析已有数据，低于 80\% 也可能值得。某些研究，例如新疗法使长期难治肺动脉高压的压力降低超过 50\%，只需两三人便足以提示值得进一步研究安全性及临床结局。

但效能不足而未发现关联的研究，难以解释与发表，宽区间会显示仍可能有较大效应。效能低却显著，也可能受批评：评审者和编辑可能怀疑是否真的预先想检验这一关联，还是检验了许多假设后，只挑显著结果。

\sourcepage{150}
\section{减少样本量并提高效能的策略}
估计人数超出现实能力时，应先检查计算，因为容易出错；再检查参数：效应是否小得不合理、变异是否过大、alpha 或 beta 是否过低、置信水平是否过高、区间是否过窄？

技术调整虽有用，但统计检验最终依赖数据中的信息。把效能从 90\% 降为 80\% 等参数改动，并不增加数据数量或质量。以下策略则能增加信息，在给定人数下提高效能，或减少所需人数。许多会改变假设，应判断新假设是否仍回答原来问题。

\subsection{使用连续变量}
连续变量通常比二分类变量节省人数。例如，以 mm Hg 记录血压，比只分高血压有无，在给定效能下人数更少，或给定人数下效能更高。
\begin{bookbox}{例 6.7\quad 连续与二分类变量的比较}
\textbf{问题：}考虑一项对照试验，研究与常规照护相比，一项新的活动计划会怎样影响养老院老年居民的满意度。既往研究发现，满意度评分（0～100 分）近似正态分布，均值为 65 分，标准差为 10 分；约 10\% 居民的满意度差或很差，定义为评分不超过 49 分。
\sourcepage{151}
如果三个月后平均满意度能提高 5 分，这项新活动计划就被认为值得开展。均值的这一变化，对应满意度差或很差者的比例降至约 5\%。一种设计把满意度作为二分类变量，即差或很差与一般及以上；另一种则使用测量中的全部信息，把满意度视为连续变量。双侧 $\alpha=0.05$、$\beta=0.20$ 时，两种设计各需要多少参与者？

\textbf{解答：二分类结局。}零假设为：养老院老年居民参加活动计划三个月后，满意度差或很差者的比例，与接受常规照护三个月后相同。备择假设为：这两个比例不同。常规照护组的预期比例 $P_0=0.10$，活动组的预期比例 $P_1=0.05$，所以效应量为 $0.10-0.05=0.05$。设定双侧 $\alpha=0.05$、$\beta=0.20$。在表 6B.1 中找到较小比例 0.05 所在行和预期差值 0.05 所在列，取对应 alpha、beta 的中间数，得出每组需要 473 人。

\textbf{连续结局。}
\sourcepage{152}
零假设为：养老院老年居民参加活动计划三个月后的平均满意度，与接受常规照护三个月后相同。备择假设为：这两个均值不同。预期常规照护组的均值为 65 分，活动组的均值为 70 分，所以效应量为 $70-65=5$ 分。满意度评分的标准差为 10 分，因此标准化效应量为 $E/S=5/10=0.5$。在双侧 $\alpha=0.05$、$\beta=0.20$ 下，表 6A 显示每组约需 64 人；前述近似式给出的 $16/(0.5)^2=64$ 人／组也相同。可见，使用连续结局变量会大幅减少所需人数，不过它回答的是不同的研究问题。
\end{bookbox}

\subsection{使用配对测量}
某些连续结局实验或队列研究，可以每人在研究开始及结束各测一次，以两次变化为结局。随后用配对测量的 $t$ 检验比较两组平均变化。每人与自身比较，去掉个体间初始差异，常可减少人数。但若两次至少没有中度相关，双次测量增加的变异就会抵消好处。临界相关为 0.5，
\sourcepage{153}
可由下式看出：
\[
\mathrm{SD}_{\text{配对}}=\mathrm{SD}_{\text{单次}}\sqrt{2(1-r_{\text{研究开始，随访}})}.
\]
例如，单次体重标准差 13 kg，两次相关为 0.8，则配对变化标准差为
\[
8.2=13\sqrt{2(1-0.8)}\quad\text{（kg）}.
\]
计算仍按通常方法（例 6.8），只是表 6A 的标准化效应量，应为预期平均变化差除以变化值标准差。
\begin{bookbox}{例 6.8\quad 配对测量的 $t$ 检验}
\textbf{问题：}回忆例 6.1：研究者希望确定，与单用沙丁胺醇相比，加用异丙托溴铵能否使 FEV$_1$ 提高 200 mL。计算显示每组需要 100 人，超出可能获得的人数。幸好，同事指出，哮喘患者即使在治疗前，FEV$_1$ 也有很大差别。这些患者之间的差异解释了治疗后 FEV$_1$ 的相当多变异，因此掩盖了治疗效应。她建议用“双样本配对 $t$ 检验”比较两组 FEV$_1$ 的变化。现有电子病历数据显示，开始新的支气管扩张治疗后，FEV$_1$ 变化值的标准差约为 250 mL。仍采用例 6.1 的效应量，即 $0.2\text{ L}=200\text{ mL}$，并设双侧 $\alpha=0.05$、$\beta=0.20$，每组需要多少人？

\textbf{解答：}
\sourcepage{154}
零假设为：治疗两周后，单用沙丁胺醇的哮喘患者与联合使用沙丁胺醇和异丙托溴铵的哮喘患者，平均 FEV$_1$ 变化相同。备择假设为：两组的平均 FEV$_1$ 变化不同。效应量 $E=200$ mL，结局变量的标准差 $S=250$ mL，所以标准化效应量 $E/S=200/250=0.8$。设定双侧 $\alpha=0.05$、$\beta=1-0.80=0.20$。表 6A 显示每组约需 26 人，远少于例 6.1 的每组 100 人。
\end{bookbox}
配对测量虽能提高效率，但可能难以获得变化标准差或计算它所需的两次相关程度。资料不足部分提供建议。
\begin{bookbox}{简短技术说明}
本章所说\textbf{双样本 $t$ 检验}，均是比较两组参与者连续变量均值。比较变量本身时称非配对（例 6.1）；比较两次测量变化时称配对（例 6.8）。\footnote{译注：本书把每个参与者两次测量变化值的组间比较称为“双样本配对 $t$ 检验”，此处沿用这一表述；配对发生在同一参与者前后两次测量，并非指两组成员彼此配对。}

第三种，\textbf{单样本配对 $t$ 检验}，比较同一组个体两次测量的平均变化是否偏离零。这常用于时间序列设计（第 12 章），即对难以随机分配的治疗采用前后比较，例如择期子宫切除对 0～10 分生活质量的影响；少有女性愿意抛硬币决定手术。
\sourcepage{155}
但没有对照组，就难以知道若不治疗会怎样。此类研究总人数为附录 6A 每组人数的一半。例如双侧 alpha 0.05、beta 0.2，检出标准差一半的变化，$E/S=0.5$，需 $64/2=32$ 人。附录 6F 进一步说明单、双样本 $t$ 检验的使用与误用。
\end{bookbox}

\subsection{使用更精密的变量}
更精密的测量能够减少变异，因此可以减少所需样本量。即使精密度只有不大的改善，也可能明显影响所需人数。例如，用 $t$ 检验估计样本量时，结局变量的标准差减少 20\%，所需人数就会减少 36\%。第 4 章介绍了重复测量等提高变量精密度的方法。

\subsection{使用大小不等的组}
总人数固定时，两组等大通常效能最高，所以附表 6A、6B.1、6B.2 都假定等大。但现实分布可能不均，一组也可能更易或更便宜招募。例如队列比较 20\% 吸烟者与 80\% 不吸烟者，或病例少、对照多的病例对照研究。一组扩大至另一组两倍，效能增加相当明显；三倍、四倍时，增益逐步减小。
\sourcepage{156}
不等组可以用软件或网站计算。

对于危险因素及结局均为二分类、每例配 $c$ 个对照的病例对照研究，有一个估计样本量的近似式（30）。在给定 alpha、beta 及效应量下，设一例配一对照所需的病例数为 $n$；每例配 $c$ 个对照时，约需病例数 $n'$、对照数 $cn'$（例 6.9）：
\[
n'=\frac{c+1}{2c}n.
\]
例如 $c=2$，只需原病例数的 $3/4$。$c$ 越大，$n'$ 越接近 $n$ 的 50\%；$c=10$ 时为 $11n/20$。
\begin{bookbox}{例 6.9\quad 每例使用多个对照}
\textbf{问题：}研究家用杀虫剂暴露是否为再生障碍性贫血危险因素。最初一例一对照需 25 例病例，但只能获得 18 例，怎么办？

\textbf{解答：}研究者应考虑每例使用多个对照，因为没有再生障碍性贫血的人很多。例如，每例配三个对照时，约需病例数为 $[(3+1)/(2\times3)]\times25=17$ 例。
\end{bookbox}

\subsection{使用更常见的结局}
二分类结局越常发生，直到约一半人发生时，效能越大。因此改变结局定义，是提高效能的好方法之一：发生越多，越有机会检出预测因素。
\sourcepage{157}
实际上，效能更依赖发生结局者人数，而非总人数。健康女性乳腺癌等罕见结局，需要很大样本。

办法之一是纳入高风险者，例如有乳腺癌家族史的女性；也可延长随访积累更多事件，或放宽定义，例如纳入导管原位癌。但这些方法（例 6.10）可能改变问题，应谨慎。
\begin{bookbox}{例 6.10\quad 使用更常见的结局}
\textbf{问题：}比较消毒漱口水与安慰剂预防上呼吸道感染。预期 200 名大学生自愿者不足，部分因三个月随访只有约 20\% 感染。如何修改？

\textbf{解答：}可以：（a）改研究感染率较高的儿科实习及住院医师；（b）在感染较常见的冬季开展；（c）随访六或十二个月；（d）在知情同意下让参与者接触鼻病毒。这些都修改了假设，但看来仍合理地服务于漱口水疗效这一总体问题。
\end{bookbox}

\section{资料不足时如何估计样本量}
缺少计算参数，常使研究者沮丧，
\sourcepage{158}
尤其用自编工具时，例如新问卷比较压力性与急迫性尿失禁女性生活质量，怎样决定评分标准差的几分之几才有临床意义？

首先，广泛搜集该主题及类似问题的既往资料。大致相似的场景、质量一般或年代较早的发现，也许足够。例如，其他泌尿系统问题，乃至结肠造口等略相关疾病的生活质量资料，是否可用？若文献无收获，应请教同行的判断，及他们是否知道相关未发表结果。

对近似钟形连续变量，忽略极端值后，常见范围的最大与最小值之差的四分之一，可近似标准差。例如，多数血钠 135～143 mEq/L，跨度 8，标准差约 2 mEq/L。

若仍无资料，可在主要研究之前开展先导研究，或获得数据集进行二次分析。几乎所有使用新工具、新测量或新招募方法的研究，都建议先导研究。它使主要研究规划实施更好，长远看节约时间；熟悉关键测量总是好事。

分类变量，或均值、标准差仍不清楚的连续变量，都可二分。把类别合为两组，或按均值、中位数或临床界点划分。例如，把生活质量分为高于中位数与不高于中位数，就免去估计标准差，
\sourcepage{159}
但仍需估计两研究组各自高于总体中位数的比例。卡方检验可给出合理、但可能略高的人数估计。

有时必须按临床意义选择效应，应请同行评议。例如，严重难治胃轻瘫自发改善最多 5\%，新侵入性疗法昂贵且副作用明显。胃肠病同行表示，若有效，最多愿治疗五人来让一人持续获益。\term{需治疗人数}（number needed to treat，NNT）为 5，对应绝对风险差 20\%，因为 NNT 为风险差倒数。因此比较 $P_0=5\%$ 与 $P_1=25\%$，在效能 0.80、双侧 alpha 0.05 下，每组需 59 人。

若全无办法，就作有根据的猜测，在几套假定下估计人数或效应量。思考问题、设想发现，常能得到合理估计，这正是规划目的。总比毫无理由便决定以双侧 alpha 0.05、80\% 效能检出标准化效应 0.5 好；顺便说，这需每组 64 人。很少资助评审者会接受那样武断的选择。

\section{应避免的常见错误}
不少新手，甚至有经验者，也会犯下列错误：
\sourcepage{160}
\begin{enumerate}
\item 设计很晚才估算；应在还能根本改变时尽早做。
\item 试验或队列中，不应假定结局与目标人群同样常见，因为参与者常比平均健康。骨质疏松性骨折研究曾估计样本髋骨折风险仅为总体三分之二，后来证实接近实际。
\item 百分比或率可使二分类变量看似连续。例如存活百分比，或每 100 人年卒中率，实际个人结局仍为生死、有无等二分类；简单估算可用卡方检验。
\item 估计的是有结局资料的人数，不是需要招募的人数；须考虑退出及缺失资料。
\item 附表假定两组等大；不等大时应使用软件或在线计算器。
\item 用 $t$ 检验时标准差很关键。结局若为变化，应使用变化的标准差，而非单次测量标准差。
\item 注意群组资料。若人数似有两个“层次”，如医师及其患者，可能存在群组问题，附表不适用。
\sourcepage{161}
\item 难以估算时，检查假设是否符合第 5 章的简单、具体、预先提出标准。
\item alpha 0.05 并不神奇，更非神圣，只是广泛采用的惯例。
\end{enumerate}

\section{小结}
\begin{enumerate}
\item 估计分析性研究的样本量时，需要完成以下步骤：陈述零假设及备择假设，明确采用单侧还是双侧；按预测变量和结局变量的类型选择检验，两者均为二分类变量时用卡方检验，一项为二分类变量而另一项为连续变量时用 $t$ 检验，两者均为连续变量时用相关系数；估计效应量，并在必要时估计测量的变异程度；根据避免 I 类及 II 类错误的重要性，设定适当的 alpha 和 beta。
\item 其他考虑包括潜在退出、分类变量、生存分析、群组、多变量调整，以及等效和非劣效试验的特殊方法。
\item 无假设的描述性研究，需要估计二分类比例或连续变量标准差，指定区间宽度及置信水平。
\item 人数预定时，可反推可检出效应，较少见的是反推效能。
\item 减少人数的方法包括连续变量、精密测量、配对测量、常见结局，以及病例对照中增加每例对照数。
\item 资料不足时，应查阅相关领域文献、咨询同行，选择有临床意义的效应量。
\item 避免估算太晚、误认百分比性质、忽视流失和缺失，以及不当处理群组或配对资料。
\end{enumerate}
\sourcepage{162}
\begin{references}
\item Lehr R. Sixteen S-squared over D-squared: a relation for crude sample size estimates. Stat Med. 1992;11:1099-1102.
\item Li H, Wang L, Wei L, Quan H. Sample size calculation for count data in comparative clinical trials with nonuniform patient accrual and early dropout. J Biopharm Stat. 2015;25(1):1-15.
\item Barthel FM, Babiker A, Royston P, Parmar MK. Evaluation of sample size and power for multi-arm survival trials allowing for non-uniform accrual, non- proportional hazards, loss to follow-up and cross-over. Stat Med. 2006;25(15):2521-2542.
\item Ahnn S, Anderson SJ. Sample size determination in complex clinical trials comparing more than two groups for survival endpoints. Stat Med. 1998;17(21):2525-2534.
\item Kerry SM, Bland JM. Trials which randomize practices II: sample size. Fam Pract. 1998;15:84-87.
\item Hemming K, Girling AJ, Sitch AJ, et al. Sample size calculations for cluster randomised controlled trials with a fixed number of clusters. BMC Med Res Methodol. 2011;11:102.
\item Jahn-Eimermacher A, Ingel K, Schneider A. Sample size in cluster-randomized trials with time to event as the primary endpoint. Stat Med. 2013;32(5):739-751.
\item Edwardes MD. Sample size requirements for case-control study designs. BMC Med Res Methodol. 2001;1:11.
\item Drescher K, Timm J, Jöckel KH. The design of case-control studies: the effect of confounding on sample size requirements. Stat Med. 1990;9:765-776.
\item Lui KJ. Sample size determination for case-control studies: the influence of the joint distribution of exposure and confounder. Stat Med. 1990;9:1485-1493.
\item Latouche A, Porcher R, Chevret S. Sample size formula for proportional hazards modelling of competing risks. Stat Med. 2004;23(21):3263-3274.
\item Novikov I, Fund N, Freedman LS. A modified approach to estimating sample size for simple logistic regression with one continuous covariate. Stat Med. 2010;29(1):97-107.
\item Vaeth M, Skovlund E. A simple approach to power and sample size calculations in logistic regression and Cox regression models. Stat Med. 2004;23(11):1781-1792.
\item Dupont WD, Plummer WD Jr. Power and sample size calculations for studies involving linear regression. Control Clin Trials. 1998;19:589-601.
\sourcepage{163}
\item Murcray CE, Lewinger JP, Conti DV, et al. Sample size requirements to detect gene-environment interactions in genome-wide association studies. Genet Epidemiol. 2011;35(3):201-210.
\item Wang S, Zhao H. Sample size needed to detect gene-gene interactions using linkage analysis. Ann Hum Genet. 2007;71(Pt 6):828-842.
\item Witte JS. Rare genetic variants and treatment response: sample size and analysis issues. Stat Med. 2012;31(25):3041-3050.
\item Willan AR. Sample size determination for cost-effectiveness trials. Pharmacoeconomics. 2011;29(11):933-949.
\item Glick HA. Sample size and power for cost-effectiveness analysis (Part 2): the effect of maximum willingness to pay. Pharmacoeconomics. 2011;29(4):287-296.
\item Glick HA. Sample size and power for cost-effectiveness analysis (Part 1). Pharmacoeconomics. 2011;29(3):189-198.
\item Patel HI. Sample size for a dose-response study. J Biopharm Stat. 1992;2:l- 8.
\item Day SJ, Graham DF. Sample size estimation for comparing two or more treatment groups in clinical trials. Stat Med. 1991;10:33-43.
\item Guo JH, Chen HJ, Luh WM. Sample size planning with the cost constraint for testing superiority and equivalence of two independent groups. Br J Math Stat Psychol. 2011;64(3):439-461.
\item Zhang P. A simple formula for sample size calculation in equivalence studies. J Biopharm Stat. 2003;13(3):529-538.
\item Stucke K, Kieser M. A general approach for sample size calculation for the three-arm ‘gold standard’ non-inferiority design. Stat Med. 2012;31(28):3579-3596.
\item Julious SA, Owen RJ. A comparison of methods for sample size estimation for non-inferiority studies with binary outcomes. Stat Methods Med Res. 2011;20(6):595-612.
\item Obuchowski NA. Sample size tables for receiver operating characteristic studies. AJR Am J Roentgenol. 2000;175(3):603-608.
\item Simel DL, Samsa GP, Matchar DB. Likelihood ratios with confidence: sample size estimation for diagnostic test studies. J Clin Epidemiol. 1991;44:763-770.
\item Sim J, Wright CC. The kappa statistic in reliability studies: use, interpretation, and sample size requirements. Phys Ther. 2005;85(3):257-268.
\item Jewell NP. Statistics for Epidemiology. Chapman and Hall; 2004:68.
\end{references}

\sourcepage{164}
\clearpage
\chapterappendix{6A}{用 \texorpdfstring{$t$}{t} 检验比较连续变量均值时，每个等大组的所需样本量}
\sourcepage{165}
\subsection*{表 6A\quad 比较两个均值的每组样本量}
最左列为标准化效应量 $E/S$。
{\small\setlength{\tabcolsep}{4pt}
\begin{longtable}{lrrrrrrrrr}
\toprule
单侧 $\alpha$ & \multicolumn{3}{c}{0.005} & \multicolumn{3}{c}{0.025} & \multicolumn{3}{c}{0.05}\\
双侧 $\alpha$ & \multicolumn{3}{c}{0.01} & \multicolumn{3}{c}{0.05} & \multicolumn{3}{c}{0.10}\\
$\beta$ & 0.05 & 0.10 & 0.20 & 0.05 & 0.10 & 0.20 & 0.05 & 0.10 & 0.20\\\midrule
\endfirsthead
\toprule
单侧 $\alpha$ & \multicolumn{3}{c}{0.005} & \multicolumn{3}{c}{0.025} & \multicolumn{3}{c}{0.05}\\
双侧 $\alpha$ & \multicolumn{3}{c}{0.01} & \multicolumn{3}{c}{0.05} & \multicolumn{3}{c}{0.10}\\
$\beta$ & 0.05 & 0.10 & 0.20 & 0.05 & 0.10 & 0.20 & 0.05 & 0.10 & 0.20\\\midrule
\endhead
0.10 & 3565 & 2978 & 2338 & 2600 & 2103 & 1571 & 2166 & 1714 & 1238\\
0.15 & 1586 & 1325 & 1040 & 1157 & 935 & 699 & 963 & 762 & 551\\
0.20 & 893 & 746 & 586 & 651 & 527 & 394 & 542 & 429 & 310\\
0.25 & 572 & 478 & 376 & 417 & 338 & 253 & 347 & 275 & 199\\
0.30 & 398 & 333 & 262 & 290 & 235 & 176 & 242 & 191 & 139\\
0.40 & 225 & 188 & 148 & 164 & 133 & 100 & 136 & 108 & 78\\
0.50 & 145 & 121 & 96 & 105 & 86 & 64 & 88 & 70 & 51\\
0.60 & 101 & 85 & 67 & 74 & 60 & 45 & 61 & 49 & 36\\
0.70 & 75 & 63 & 50 & 55 & 44 & 34 & 45 & 36 & 26\\
0.80 & 58 & 49 & 39 & 42 & 34 & 26 & 35 & 28 & 21\\
0.90 & 46 & 39 & 32 & 34 & 27 & 21 & 28 & 22 & 16\\
1.00 & 38 & 32 & 26 & 27 & 23 & 17 & 23 & 18 & 14\\\bottomrule
\end{longtable}}

\sourcepage{166}
\noindent\small $E/S$ 为预期效应量 $E$（两组均值差）除以结局标准差 $S$。按该行与指定 alpha、beta 列查得每组人数。单样本 $t$ 检验的总人数，为表列数字的一半。\par\normalsize
\subsection{计算变异程度}
变异常报告为标准差，或均值标准误（standard error of the mean，SEM）。样本量计算最有用的是标准差，转换很容易：标准差等于标准误乘以 $\sqrt N$，其中 $N$ 为用于计算平均值的参与者人数。例如，25 人低纤维饮食后减重为 $10\pm2$ kg（均值 $\pm$ SEM），则标准差为 $2\text{ kg}\times\sqrt{25}=10\text{ kg}$。
\subsection{其他取值}
表中没有的 $E/S$，或两组大小不等时，请使用 \url{https://www.sample-size.net} 计算器。

\clearpage
\chapterappendix{6B}{用卡方统计量或 \texorpdfstring{$Z$}{Z} 检验比较二分类变量比例时，每个等大组的所需样本量}
\sourcepage{167}
\subsection*{表 6B.1\quad 比较两个比例的每组样本量}
每格三数分别对应：上方，单侧 $\alpha=0.05$ 或双侧 $\alpha=0.10$，$\beta=0.20$；中间，单侧 $\alpha=0.025$ 或双侧 $\alpha=0.05$，$\beta=0.20$；下方，单侧 $\alpha=0.025$ 或双侧 $\alpha=0.05$，$\beta=0.10$。各数值列为 $|P_1-P_0|$，最左列为两比例中较小者。
{\small\setlength{\tabcolsep}{3.5pt}
\begin{longtable}{lrrrrrrrrrr}
\toprule $\min(P_0,P_1)$ & 0.05 & 0.10 & 0.15 & 0.20 & 0.25 & 0.30 & 0.35 & 0.40 & 0.45 & 0.50\\\midrule
\endfirsthead
\toprule $\min(P_0,P_1)$ & 0.05 & 0.10 & 0.15 & 0.20 & 0.25 & 0.30 & 0.35 & 0.40 & 0.45 & 0.50\\\midrule
\endhead
0.05 & 381 & 129 & 72 & 47 & 35 & 27 & 22 & 18 & 15 & 13\\*
 & 473 & 159 & 88 & 59 & 43 & 33 & 26 & 22 & 18 & 16\\*
 & 620 & 207 & 113 & 75 & 54 & 41 & 33 & 27 & 23 & 19\\\midrule
0.10 & 578 & 175 & 91 & 58 & 41 & 31 & 24 & 20 & 16 & 14\\*
 & 724 & 219 & 112 & 72 & 51 & 37 & 29 & 24 & 20 & 17\\*
 & 958 & 286 & 146 & 92 & 65 & 48 & 37 & 30 & 25 & 21\\\midrule
0.15 & 751 & 217 & 108 & 67 & 46 & 34 & 26 & 21 & 17 & 15\\*
 & 944 & 270 & 133 & 82 & 57 & 41 & 32 & 26 & 21 & 18\\*
 & 1252 & 354 & 174 & 106 & 73 & 53 & 42 & 33 & 26 & 22\\\midrule
0.20 & 900 & 251 & 121 & 74 & 50 & 36 & 28 & 22 & 18 & 15\\*
 & 1133 & 313 & 151 & 91 & 62 & 44 & 34 & 27 & 22 & 18\\*
 & 1504 & 412 & 197 & 118 & 80 & 57 & 44 & 34 & 27 & 23\\\midrule
0.25 & 1024 & 278 & 132 & 79 & 53 & 38 & 29 & 23 & 18 & 15\\*
 & 1289 & 348 & 165 & 98 & 66 & 47 & 35 & 28 & 22 & 18\\*
 & 1714 & 459 & 216 & 127 & 85 & 60 & 46 & 35 & 28 & 23\\\midrule
0.30 & 1123 & 300 & 141 & 83 & 55 & 39 & 29 & 23 & 18 & 15\\*
 & 1415 & 376 & 175 & 103 & 68 & 48 & 36 & 28 & 22 & 18\\*
 & 1883 & 496 & 230 & 134 & 88 & 62 & 47 & 36 & 28 & 23\\\midrule
0.35 & 1197 & 315 & 146 & 85 & 56 & 39 & 29 & 23 & 18 & 15\\*
 & 1509 & 395 & 182 & 106 & 69 & 48 & 36 & 28 & 22 & 18\\*
 & 2009 & 522 & 239 & 138 & 90 & 62 & 47 & 35 & 27 & 22\\\midrule
0.40 & 1246 & 325 & 149 & 86 & 56 & 39 & 29 & 22 & 17 & 14\\*
 & 1572 & 407 & 186 & 107 & 69 & 48 & 35 & 27 & 21 & 17\\*
 & 2093 & 538 & 244 & 139 & 90 & 62 & 46 & 34 & 26 & 21\\\midrule
0.45 & 1271 & 328 & 149 & 85 & 55 & 38 & 28 & 21 & 16 & 13\\*
 & 1603 & 411 & 186 & 106 & 68 & 47 & 34 & 26 & 20 & 16\\*
 & 2135 & 543 & 244 & 138 & 88 & 60 & 44 & 33 & 25 & 19\\\midrule
0.50 & 1271 & 325 & 146 & 83 & 53 & 36 & 26 & 20 & 15 & —\\*
 & 1603 & 407 & 182 & 103 & 66 & 44 & 32 & 24 & 18 & —\\*
 & 2135 & 538 & 239 & 134 & 85 & 57 & 42 & 30 & 23 & —\\\midrule
0.55 & 1246 & 315 & 141 & 79 & 50 & 34 & 24 & 18 & — & —\\*
 & 1572 & 395 & 175 & 98 & 62 & 41 & 29 & 22 & — & —\\*
 & 2093 & 522 & 230 & 127 & 80 & 53 & 37 & 27 & — & —\\\midrule
0.60 & 1197 & 300 & 132 & 74 & 46 & 31 & 22 & — & — & —\\*
 & 1509 & 376 & 165 & 91 & 57 & 37 & 26 & — & — & —\\*
 & 2009 & 496 & 216 & 118 & 73 & 48 & 33 & — & — & —\\\midrule
0.65 & 1123 & 278 & 121 & 67 & 41 & 27 & — & — & — & —\\*
 & 1415 & 348 & 151 & 82 & 51 & 33 & — & — & — & —\\*
 & 1883 & 459 & 197 & 106 & 65 & 41 & — & — & — & —\\\midrule
0.70 & 1024 & 251 & 108 & 58 & 35 & — & — & — & — & —\\*
 & 1289 & 313 & 133 & 72 & 43 & — & — & — & — & —\\*
 & 1714 & 412 & 174 & 92 & 54 & — & — & — & — & —\\\midrule
0.75 & 900 & 217 & 91 & 47 & — & — & — & — & — & —\\*
 & 1133 & 270 & 112 & 59 & — & — & — & — & — & —\\*
 & 1504 & 354 & 146 & 75 & — & — & — & — & — & —\\\midrule
0.80 & 751 & 175 & 72 & — & — & — & — & — & — & —\\*
 & 944 & 219 & 88 & — & — & — & — & — & — & —\\*
 & 1252 & 286 & 113 & — & — & — & — & — & — & —\\\midrule
0.85 & 578 & 129 & — & — & — & — & — & — & — & —\\*
 & 724 & 159 & — & — & — & — & — & — & — & —\\*
 & 958 & 207 & — & — & — & — & — & — & — & —\\\midrule
0.90 & 381 & — & — & — & — & — & — & — & — & —\\*
 & 473 & — & — & — & — & — & — & — & — & —\\*
 & 620 & — & — & — & — & — & — & — & — & —\\\midrule
\end{longtable}}

\sourcepage{168}
单侧估计使用 $Z$ 统计量。$P_0$ 表示一组中预期发生结局的比例，$P_1$ 是另一组的比例。在病例对照研究，$P_1$ 是病例中具有预测特征者比例，$P_0$ 是对照中相应比例。按较小的 $P_0$ 或 $P_1$ 行和预期绝对差列读取；三个数表示指定 alpha、beta 下各组所需样本量。$P_0$ 或 $P_1$ 为 0.01～0.10 的更详细数字见表 6B.2。
\subsection{其他取值}
\sourcepage{169}
两表未列的比例，或两组大小不等时，请用网站计算器。
\subsection*{表 6B.2\quad 比较两个比例的每组样本量：较小比例为 0.01～0.10}
上、中、下三数的 alpha、beta 组合与表 6B.1 相同；各列为两比例绝对差，最左列为较小比例。
{\small\setlength{\tabcolsep}{3.5pt}
\begin{longtable}{lrrrrrrrrrr}
\toprule $\min(P_0,P_1)$ & 0.01 & 0.02 & 0.03 & 0.04 & 0.05 & 0.06 & 0.07 & 0.08 & 0.09 & 0.10\\\midrule
\endfirsthead
\toprule $\min(P_0,P_1)$ & 0.01 & 0.02 & 0.03 & 0.04 & 0.05 & 0.06 & 0.07 & 0.08 & 0.09 & 0.10\\\midrule
\endhead
0.01 & 2019 & 700 & 396 & 271 & 204 & 162 & 134 & 114 & 98 & 87\\*
 & 2512 & 864 & 487 & 332 & 249 & 197 & 163 & 138 & 120 & 106\\*
 & 3300 & 1125 & 631 & 428 & 320 & 254 & 209 & 178 & 154 & 135\\\midrule
0.02 & 3205 & 994 & 526 & 343 & 249 & 193 & 157 & 131 & 113 & 97\\*
 & 4018 & 1237 & 651 & 423 & 306 & 238 & 192 & 161 & 137 & 120\\*
 & 5320 & 1625 & 852 & 550 & 397 & 307 & 248 & 207 & 177 & 154\\\midrule
0.03 & 4367 & 1283 & 653 & 414 & 294 & 224 & 179 & 148 & 126 & 109\\*
 & 5493 & 1602 & 813 & 512 & 363 & 276 & 220 & 182 & 154 & 133\\*
 & 7296 & 2114 & 1067 & 671 & 474 & 359 & 286 & 236 & 199 & 172\\\midrule
0.04 & 5505 & 1564 & 777 & 482 & 337 & 254 & 201 & 165 & 139 & 119\\*
 & 6935 & 1959 & 969 & 600 & 419 & 314 & 248 & 203 & 170 & 146\\*
 & 9230 & 2593 & 1277 & 788 & 548 & 410 & 323 & 264 & 221 & 189\\\midrule
0.05 & 6616 & 1838 & 898 & 549 & 380 & 283 & 222 & 181 & 151 & 129\\*
 & 8347 & 2308 & 1123 & 686 & 473 & 351 & 275 & 223 & 186 & 159\\*
 & 11123 & 3061 & 1482 & 902 & 620 & 460 & 360 & 291 & 242 & 206\\\midrule
0.06 & 7703 & 2107 & 1016 & 615 & 422 & 312 & 243 & 197 & 163 & 139\\*
 & 9726 & 2650 & 1272 & 769 & 526 & 388 & 301 & 243 & 202 & 171\\*
 & 12973 & 3518 & 1684 & 1014 & 691 & 508 & 395 & 318 & 263 & 223\\\midrule
0.07 & 8765 & 2369 & 1131 & 680 & 463 & 340 & 263 & 212 & 175 & 148\\*
 & 11076 & 2983 & 1419 & 850 & 577 & 423 & 327 & 263 & 217 & 183\\*
 & 14780 & 3965 & 1880 & 1123 & 760 & 555 & 429 & 343 & 283 & 239\\\midrule
0.08 & 9803 & 2627 & 1244 & 743 & 502 & 367 & 282 & 227 & 187 & 158\\*
 & 12393 & 3308 & 1562 & 930 & 627 & 457 & 352 & 282 & 232 & 195\\*
 & 16546 & 4401 & 2072 & 1229 & 827 & 602 & 463 & 369 & 303 & 255\\\midrule
0.09 & 10816 & 2877 & 1354 & 804 & 541 & 393 & 302 & 241 & 198 & 167\\*
 & 13679 & 3626 & 1702 & 1007 & 676 & 491 & 377 & 300 & 246 & 207\\*
 & 18270 & 4827 & 2259 & 1333 & 893 & 647 & 495 & 393 & 322 & 270\\\midrule
0.10 & 11804 & 3121 & 1461 & 863 & 578 & 419 & 320 & 255 & 209 & 175\\*
 & 14933 & 3936 & 1838 & 1083 & 724 & 523 & 401 & 318 & 260 & 218\\*
 & 19952 & 5242 & 2441 & 1434 & 957 & 690 & 527 & 417 & 341 & 285\\\midrule
\end{longtable}}

\sourcepage{170}
单侧估计使用 $Z$ 统计量。
\subsection{其他取值}
两表未列的 $P_0$、$P_1$，或两组大小不等时，请使用 \url{https://www.sample-size.net}。

\clearpage
\chapterappendix{6C}{使用相关系数时所需的总样本量}
\sourcepage{171}
\subsection*{表 6C\quad 判断相关系数是否异于零的样本量}
最左列为预期相关系数 $r$。按该行和指定 alpha、beta 列查总人数。
{\small\setlength{\tabcolsep}{4pt}
\begin{longtable}{lrrrrrrrrr}
\toprule
单侧 $\alpha$ & \multicolumn{3}{c}{0.005} & \multicolumn{3}{c}{0.025} & \multicolumn{3}{c}{0.05}\\
双侧 $\alpha$ & \multicolumn{3}{c}{0.01} & \multicolumn{3}{c}{0.05} & \multicolumn{3}{c}{0.10}\\
$\beta$ & 0.05 & 0.10 & 0.20 & 0.05 & 0.10 & 0.20 & 0.05 & 0.10 & 0.20\\\midrule
\endfirsthead
\toprule
单侧 $\alpha$ & \multicolumn{3}{c}{0.005} & \multicolumn{3}{c}{0.025} & \multicolumn{3}{c}{0.05}\\
双侧 $\alpha$ & \multicolumn{3}{c}{0.01} & \multicolumn{3}{c}{0.05} & \multicolumn{3}{c}{0.10}\\
$\beta$ & 0.05 & 0.10 & 0.20 & 0.05 & 0.10 & 0.20 & 0.05 & 0.10 & 0.20\\\midrule
\endhead
0.05 & 7118 & 5947 & 4663 & 5193 & 4200 & 3134 & 4325 & 3424 & 2469\\
0.10 & 1773 & 1481 & 1162 & 1294 & 1047 & 782 & 1078 & 854 & 616\\
0.15 & 783 & 655 & 514 & 572 & 463 & 346 & 477 & 378 & 273\\
0.20 & 436 & 365 & 287 & 319 & 259 & 194 & 266 & 211 & 153\\
0.25 & 276 & 231 & 182 & 202 & 164 & 123 & 169 & 134 & 98\\
0.30 & 189 & 158 & 125 & 139 & 113 & 85 & 116 & 92 & 67\\
0.35 & 136 & 114 & 90 & 100 & 82 & 62 & 84 & 67 & 49\\
0.40 & 102 & 86 & 68 & 75 & 62 & 47 & 63 & 51 & 37\\
0.45 & 79 & 66 & 53 & 58 & 48 & 36 & 49 & 39 & 29\\
0.50 & 62 & 52 & 42 & 46 & 38 & 29 & 39 & 31 & 23\\
0.60 & 40 & 34 & 27 & 30 & 25 & 19 & 26 & 21 & 16\\
0.70 & 27 & 23 & 19 & 20 & 17 & 13 & 17 & 14 & 11\\
0.80 & 18 & 15 & 13 & 14 & 12 & 9 & 12 & 10 & 8\\\bottomrule
\end{longtable}}

\subsection{其他取值}
表中未列的 $r$，请使用网站计算器。

\sourcepage{172}
\clearpage
\chapterappendix{6D}{连续变量描述性研究的样本量}
\sourcepage{173}
\subsection*{表 6D\quad 常见 \texorpdfstring{$W/S$}{W/S} 下的样本量}
数字列为置信水平。
{\small\setlength{\tabcolsep}{4pt}
\begin{longtable}{lrrr}
\toprule $W/S$ & 90\% & 95\% & 99\%\\\midrule
\endfirsthead
\toprule $W/S$ & 90\% & 95\% & 99\%\\\midrule
\endhead
0.10 & 1083 & 1537 & 2665\\
0.15 & 482 & 683 & 1180\\
0.20 & 271 & 385 & 664\\
0.25 & 174 & 246 & 425\\
0.30 & 121 & 171 & 295\\
0.35 & 89 & 126 & 217\\
0.40 & 68 & 97 & 166\\
0.50 & 44 & 62 & 107\\
0.60 & 31 & 43 & 74\\
0.70 & 23 & 32 & 55\\
0.80 & 17 & 25 & 42\\
0.90 & 14 & 19 & 33\\
1.00 & 11 & 16 & 27\\\bottomrule
\end{longtable}}

$W/S$ 是标准化置信区间宽度，即所需总宽 $W$ 除以变量标准差 $S$。按标准化宽度行与指定置信水平列查总人数。
\subsection{其他取值}
表中未列的 $W/S$，请使用网站计算器。

\clearpage
\chapterappendix{6E}{二分类变量描述性研究的样本量}
\sourcepage{174}
\subsection*{表 6E\quad 比例的样本量}
最左列为预期比例 $P$，各数值列为置信区间总宽 $W$。每格上、中、下三数依次为 90\%、95\%、99\% 置信水平所需样本量。
{\small\setlength{\tabcolsep}{3.5pt}
\begin{longtable}{lrrrrrrr}
\toprule $P$ & 0.10 & 0.15 & 0.20 & 0.25 & 0.30 & 0.35 & 0.40\\\midrule
\endfirsthead
\toprule $P$ & 0.10 & 0.15 & 0.20 & 0.25 & 0.30 & 0.35 & 0.40\\\midrule
\endhead
0.10 & 98 & 44 & — & — & — & — & —\\*
 & 138 & 61 & — & — & — & — & —\\*
 & 239 & 106 & — & — & — & — & —\\\midrule
0.15 & 139 & 62 & 35 & 22 & — & — & —\\*
 & 196 & 87 & 49 & 31 & — & — & —\\*
 & 339 & 151 & 85 & 54 & — & — & —\\\midrule
0.20 & 174 & 77 & 44 & 28 & 19 & 14 & —\\*
 & 246 & 109 & 61 & 39 & 27 & 20 & —\\*
 & 426 & 189 & 107 & 68 & 47 & 35 & —\\\midrule
0.25 & 204 & 91 & 51 & 33 & 23 & 17 & 13\\*
 & 288 & 128 & 72 & 46 & 32 & 24 & 18\\*
 & 499 & 222 & 125 & 80 & 55 & 41 & 31\\\midrule
0.30 & 229 & 102 & 57 & 37 & 25 & 19 & 14\\*
 & 323 & 143 & 81 & 52 & 36 & 26 & 20\\*
 & 559 & 249 & 140 & 89 & 62 & 46 & 35\\\midrule
0.40 & 261 & 116 & 65 & 42 & 29 & 21 & 16\\*
 & 369 & 164 & 92 & 59 & 41 & 30 & 23\\*
 & 639 & 284 & 160 & 102 & 71 & 52 & 40\\\midrule
0.50 & 272 & 121 & 68 & 44 & 30 & 22 & 17\\*
 & 384 & 171 & 96 & 61 & 43 & 31 & 24\\*
 & 666 & 296 & 166 & 107 & 74 & 54 & 42\\\midrule
\end{longtable}}

\sourcepage{175}
按具有目标特征的预期比例行和所需区间总宽列查人数；三数对应三个置信水平。
\subsection{其他取值}
表中未列的 $P$，请使用网站计算器。

\clearpage
\chapterappendix{6F}{\texorpdfstring{$t$}{t} 检验的使用与误用}
\sourcepage{176}
本章重点之一的双样本 $t$ 检验，比较两组变量均值。组别可按二分类预测变量定义，例如随机试验的有效药与安慰剂、队列的危险因素有无；也可按二分类结局定义，例如病例对照研究。比较一次测量时为非配对；比较每人两个时间点，例如干预前后测量的变化时，本书称配对。第三种单样本配对 $t$ 检验，则比较单组个体两次测量的平均变化与零或其他指定变化值。

表 6F 展示了原本应作组间比较，却误用单样本配对检验的情况：随机盲法试验研究新安眠药对 0～10 分生活质量的影响。一些研究者曾分别对治疗及安慰剂组各做一次单样本检验，甚至发表，这并不恰当。

标有匕首号 $\dagger$ 的 $P$ 值来自单样本配对检验。治疗组 $P=0.05$，显示研究期间显著变化；对照组 $P=0.16$，无显著变化。但这不能推断组间差别，不能据此说治疗有显著效果。

星号 $*$ 是适当的双样本检验。前两个 $P=0.87$ 和 0.64 为非配对，显示初始及末次生活质量组间均无显著差异。最后 $P=0.17$ 是本书所称双样本配对检验，比末次比较 0.64 更接近 0.05，因为变化均值的标准差更小。
\sourcepage{177}
但是治疗组平均改善 1.3 分，与安慰剂 0.9 分无显著差异。因此，正确结论是本研究未发现治疗有效。
\begin{bookbox}{表 6F\quad 分析配对资料的正确与错误方式}
\begin{tabularx}{\linewidth}{YYYY}
\toprule 测量时间 & 治疗（$N=100$） & 对照（$N=100$） & $P$ 值\\
\multicolumn{4}{c}{生活质量：均值 $\pm$ 标准差}\\\midrule
研究开始 & $7.0\pm4.5$ & $7.1\pm4.4$ & $0.87^*$\\
研究结束 & $8.3\pm4.7$ & $8.0\pm4.6$ & $0.64^*$\\
变化（结束减开始） & $1.3\pm2.1$ & $0.9\pm2.0$ & $0.17^*$\\
$P$ 值 & $0.05^\dagger$ & $0.16^\dagger$ & \\\bottomrule
\end{tabularx}
\footnotesize $*$ 比较治疗组与对照组。$\dagger$ 比较变化值与零。
\end{bookbox}

\clearpage
\chapterappendix{6G}{第 6 章习题：估计样本量与效能}
\sourcepage{178}
\begin{enumerate}
\item 回顾第 5 章习题 1。研究需多少胃癌病例？若效能要求 0.90 呢？若显著性水平 0.01 呢？附加题：如果只有 60 例，她可以怎样处理？
\item 肌力随年龄增长下降。初步资料提示，部分可能由脱氢表雄酮（dehydroepiandrosterone，DHEA）逐渐缺乏造成。研究者计划随机给予老年人 DHEA 或相同外观安慰剂六个月，再测肌力。既往平均握力 20 kg、标准差 8 kg。双侧 alpha 0.05、beta 0.10 时，为显示治疗组与安慰剂至少 10\% 差异，需多少人？若 beta 0.20 呢？
\item 习题 2 计算人数超过可招募数。同事指出，老年人握力差异很大，可能掩盖治疗效果；建议研究开始及治疗后各测一次，以变化作结局。小型先导研究显示六个月变化标准差仅 2 kg。双侧 alpha 0.05、beta 0.10，仍检出效应量 2 kg，每组需多少人？
\item 研究者怀疑阅读障碍三年级学生中左利手更多。既往约 10\% 为左利手，阅读障碍不常见。计划选学区全部阅读障碍者为病例，随机选同样多无障碍者为对照。若检出左利手相对右利手的阅读障碍 OR 为 2.0，双侧 alpha 0.05、beta 0.20，需要多少人？
\item 研究者希望估计所在机构医学生平均智商，99\% 置信区间为 $\pm3$ 分。小型先导研究显示约 110～150 分，需要的大致人数是多少？
\end{enumerate}

